\documentclass[11pt]{article}
\usepackage[margin=2.3cm]{geometry}
\usepackage{amsmath,amssymb,amsthm}
\usepackage[T1]{fontenc}
\usepackage{enumitem}
\usepackage{hyperref}
\newcommand{\tri}{\mathbin{\triangleleft}}
\newcommand{\N}{\mathbb{N}}
\newcommand{\WIPBASE}{ada7c40}
\setlength{\parskip}{4pt}\setlength{\parindent}{0pt}

\title{Review: \emph{A uniqueness theorem for the container derivative}\\[2pt]
\large \S6 solidity converse and the $\N\cdot y$ sharpness claim}
\author{Rick\\[2pt]\small To: MacBeth (Kodamai)}
\date{2026-10-05}

\begin{document}
\maketitle
\thispagestyle{empty}
\begin{center}
\begin{tabular}{ll}
Reviewed: & MacBeth WIP \texttt{57bac18} (11 pp., dated 4 Oct 2026)\\
WIP base commit: & \texttt{\WIPBASE} (grandpa-rick/work-in-progress, HEAD before this review)\\
Check script: & $<$1\,s Python enumeration, reproduced in Appendix A
\end{tabular}
\end{center}

\section*{Verdict}
\textbf{Error, not a gap.} The arithmetic of \S6 is fine. The structure it is applied to is not
the one named in the theorem. Lemma~11 computes $D_2$ and $P'$ in the Weil-algebra model
$\N[\varepsilon_1,\varepsilon_2]/(\dots)$, with $\varepsilon$-truncation and ``$0\cdot x=0$''
absorption, but asserts that the computation takes place in $(\mathrm{Poly},\tri,y)$. Computed
honestly in Poly, $D_2=(1+n)^2\cdot y$ and $P'=(1+n)\cdot y$. Universality then gives
$(1+n)^2=1+n$, which kills $n=1$, and with it $\partial$. The underlying cause is upstream
(Finding~1): $T(p)=p+y\,p'$ is \emph{not} of the form $(-)\tri D$ for any polynomial $D$. So
Theorem~1 is false as stated. I do not accept \S6. The sharpness paragraph also overclaims
(Finding~4).

\section*{Findings}
\begin{enumerate}[label=\textbf{F\arabic*.},leftmargin=*]

\item \textbf{$\partial$ is not representable in $(\mathrm{Poly},\tri,y)$.} (p.\,4, ll.\,178--183;
Thm~1 p.\,2; Thm~16 p.\,9.) If $T=(-)\tri D$ then $T(y)=y\tri D=D$. For the dual-number
structure $T(y)=y+y\cdot 1=2y$, so this forces $D=2y$. But
\[
y^2\tri 2y=(2y)^2=4y^2\neq 3y^2=y^2+y\cdot(y^2)' ,\qquad y^3\tri 2y=8y^3\neq 4y^3 .
\]
Substitution has no ``mod $\varepsilon^2$''. The polynomial $y+\varepsilon v$ is just $2y$, and
$p(2y)$ keeps every cross term. The left-hand convention fails as well: $D\tri p=2p$. The
sentence ``$(-)\tri D$ is exactly the dual-number substitution $T$ of [11]'' (p.\,4, l.\,179) is
the false step. The dual-number $T$ is base change along $\N\to\N[\varepsilon]/\varepsilon^2$,
which is a tensoring in a category of $\N$-semimodules or $\N$-algebras, not a right
$\tri$-tensoring in Poly.

\item \textbf{Lemma 11 computes in the wrong category.} (p.\,7, ll.\,350--378.) The words ``We
work in the based setting: $\varepsilon$-grading modulo $\varepsilon^2$ \dots\ Weil-algebra
model'' (l.\,351) silently replace Poly by $\N$-algebras. The three counts mix the two settings:
\begin{itemize}[leftmargin=1.2em]
\item $D\tri D$: $(1+n)^2$. This is correct in both settings (true on the linear layer).
\item $D_2=D\times_y D$: the paper gets $1+2n$. In Poly, a pullback of linear polynomials over $y$
has position set $S_1\times_{1}S_2$, and its direction set is the pushout $1\leftarrow1\to1=1$.
So $D_2=(1+n)^2\cdot y$. Over a one-shape base the fibre product \emph{is} the product. The
sentence ``creates no product of the two infinitesimals'' (l.\,366) holds only in the Weil
algebra, where $\varepsilon_1\varepsilon_2$ is killed by the quotient.
\item $P'$: the paper gets $1+n+n^2$. In Poly a morphism of linear polynomials is a function on
shapes. $\hat p$ sends every shape of $D$ to the single shape of $y$, so
$\hat p\tri D:(a,b)\mapsto b$ on positions of $D\tri D$. The preimage of $\operatorname{im}\hat
z=\{0\}$ is $\{(a,0)\}$, so $P'=(1+n)\cdot y$. With the other convention, $D\tri\hat p$, the
count is the same. ``Kills the mixed term'' (l.\,370) uses the absorption
$0\cdot\varepsilon_{\rm in}=0$, and Poly has no such absorption.
\end{itemize}
The enumeration in Appendix~A gives, for $n=0,\dots,4$: Poly
$(|D_2|,|P'|)=$ $(1,1)$, $(4,2)$, $(9,3)$, $(16,4)$, $(25,5)$;
Weil $(1,1)$, $(3,3)$, $(5,7)$, $(7,13)$, $(9,21)$.
Universality holds in Poly only at $n=0$, and in the Weil model only at $n\in\{0,1\}$. This is
consistent with F1: $(-)\tri 2y$ fails the universality axiom, so it is not a tangent structure
at all.

\item \textbf{The rest of \S6 is sound \emph{conditionally}.} (p.\,8, ll.\,381--398.) This
answers your question~(a) directly.
\begin{itemize}[leftmargin=1.2em]
\item The $D$ is fixed, so $n$ is fixed. Evaluating the natural iso $T_2\cong V$ at the single
object $y$ gives one equation for that $n$. That is all you need and all you use, so there is no
``for all $n$'' issue.
\item Cancellation in $\N$ is correct: $1+2n=1+n+n^2\Rightarrow n=n^2\Rightarrow n\in\{0,1\}$.
\item Lemma~10's reduction $T_2=(-)\tri D_2$, $V=(-)\tri P'$ is plausible for $D$ linear, and I
did not find a problem with it.
\end{itemize}
So the converse is correct for the Weil-model representable structures. But those are exactly
Lanfranchi--Lemay's setting with $R=\N$ and $D=W$ an $\N$-algebra, not Poly. The ``one new idea''
(``ranks add, never subtract'') then has nothing Poly-specific left in it. The additive count
$1+n+n^2$ is just the $\N$-basis count of $\ker(W\otimes W\to W)$-over-zero, and it works over
any rig with a free $W$.

\item \textbf{Sharpness: Prop.~9 shows the method fails, not that the theorem fails.}
(p.\,6, ll.\,293--310; abstract l.\,26; Rem.~15 p.\,8; \S9 l.\,502.) This answers question~(b).
\begin{enumerate}[label=(\roman*),leftmargin=1.5em]
\item ``Consequently unqualified uniqueness over all of Poly is false'' (l.\,297) is a non
sequitur. A solid part is not a tangent structure. Your forward direction (\S7) is only
established for $M\in\{0,y\}$, and you say it ``never has to manufacture a third structure''.
To claim falsity you must exhibit $(-)\tri\N y$ with all its tangent axioms. As it stands, the
correct statement is: \emph{the rank argument does not exclude $\N\cdot y$}.
\item Solidity is about the structure map $\mu$, not about the existence of \emph{a} bijection
$\N\times\N\cong\N$. Cantor pairing shows only that the ranks admit an iso. The 40$\times$40 /
12$\times$12 ``injectivity windows'' add nothing, because $|\N\times\N|=|\N|$ is standard.
\item $\N\cdot y\notin\mathrm{SPoly}$, and $(-)\tri\N y$ does not even preserve SPoly. For example,
$y\tri\N y=\N y$. So the hypothesis is excluded by the ambient category, not by a rank argument.
Please say so. Also separate ``finite support'' (finitely many shapes) from ``finitary'' (finite
position sets). $\N\cdot y$ is finitary but not finite-support.
\item An interesting by-product of F2: in honest Poly the universality ranks are
$|D_2|=\kappa^2=\kappa=|P'|$ for every infinite $\kappa$. So infinite rank passes the test that
kills $n=1$. Remark~15's ``$1+2\kappa=\kappa=1+\kappa+\kappa^2$'' is correct cardinal arithmetic
but in the wrong model.
\end{enumerate}

\item \textbf{Linearity of $M$ is asserted, not proved.} (p.\,4, ll.\,198--203.) Poly is not
pointed: $0$ is initial and $1$ is terminal. So ``$M=\ker p$'' and the splitting $D=y\oplus M$ are
not available as written. The argument ``$\varepsilon y^2$ has no vector-space fibre'' is
heuristic. In fact the additive-bundle axiom alone does not bound $n$: any finite commutative
monoid on $1+n$ with unit the base shape gives $D_2\to D$. Please either prove linearity from the
axioms or make it a hypothesis. The Lean boundary statement (p.\,9) already lists it as a ``cited
framework fact''. Cited from where?

\item \textbf{Minor.}
\begin{itemize}[leftmargin=1.2em]
\item In Rem.~14 the claim ``solidity coincides with indecomposability'' is unsupported. $T_r$
for $D=y+\varepsilon v$ has $D_r$ of Weil rank $1+r$, which is not the object $y+\sum\varepsilon_i
v_i$ viewed as a standalone $D$.
\item The Lean certificate (pp.\,6, 9) certifies $1+2n=1+n+n^2\iff n\le1$. That is correct and
irrelevant to F1 and F2, because the error is in which numbers go into the equation. Please do not
cite it as evidence for \S6.
\end{itemize}
\end{enumerate}

\section*{Suggested fixes}
\begin{enumerate}[leftmargin=*]
\item \textbf{Change the ambient category.} Restate the theorem for representable tangent
structures $T=(-)\otimes W$ on $\N$-semimodule / $\N$-algebra objects, i.e.\ on
\emph{coefficients} of polynomials (polynomials with coefficients in a rig $R$, with $T$ given by
base change $R\to R\otimes W$). That is where your $D_2$, $P'$ counts are true, and \S6 then goes
through as written. Expect a referee to ask what is new over Lanfranchi--Lemay with $R=\N$. The
honest answer is ``free modules over a rig, no PID'', which is real but small. Drop ``(Poly,
$\tri$, $y$)'' throughout.
\item \textbf{Or keep Poly and report the true answer.} If you insist on right
$\tri$-tensoring in Poly, the finite-rank classification by your own method is
\emph{$\{\mathrm{id}\}$ only}: $(1+n)^2=1+n\iff n=0$. $\partial$ is then a non-representable
tangent structure. That is a publishable-sounding sentence in its own right, but it is the
opposite of Theorem~1.
\item \textbf{Rephrase sharpness} as ``the rank arguments of \S\S4 and 6 do not exclude
$M=\N\cdot y$; whether $(-)\tri\N y$ (or its Weil analogue) underlies a tangent structure is
open.'' Do this unless you construct one. Delete ``unqualified uniqueness over Poly is false'' from
the abstract, l.\,297 and \S9.
\item Prove or hypothesise linearity of $M$ (F5).
\item Fix the companion note [11] at the same point (``$(-)\tri D$ is exactly the dual-number
substitution'') if it makes the same identification.
\end{enumerate}

\appendix
\section{Check script (summary)}
\small
$T(p)=p+yp'$ vs.\ $p\tri 2y$: they agree on $1$ and $y$; on $y^2$, $y^3$, $y^2+1$ they give
$3y^2$ vs.\ $4y^2$, $4y^3$ vs.\ $8y^3$, $3y^2+1$ vs.\ $4y^2+1$.
Linear pullbacks are computed as fibre products of shape sets. $P'=\{(a,b)\in D\tri D: b=0\}$,
with the same count for $a=0$.
\begin{center}\begin{tabular}{c|ccc|ccc}
$n$ & Poly $|D_2|$ & Poly $|P'|$ & iso? & Weil $T_2$ & Weil $V$ & iso?\\\hline
0&1&1&yes&1&1&yes\\
1&4&2&no&3&3&yes\\
2&9&3&no&5&7&no\\
3&16&4&no&7&13&no\\
4&25&5&no&9&21&no
\end{tabular}\end{center}

\end{document}
